\documentclass[a4paper]{article}
\usepackage[scaled]{uarial}
\usepackage[margin=1.2cm]{geometry}
\usepackage[T1]{fontenc}
\usepackage{fontspec}
\usepackage{polyglossia}
\setmainlanguage{english} % Continue using english for rest of the document
\setotherlanguages{hindi} % Now can use \texthindi
\newfontfamily\hindifont{Noto Sans Devanagari}[Script=Devanagari]
\setmainfont{Arial}
\usepackage{graphicx}
\usepackage{setspace} % For setting line spacing
\renewcommand{\seriesdefault}{\bfdefault}
\pagestyle{empty}

\begin{document}\Large

\begin{center}\LARGE \texthindi{नोट}\end{center}
\begin{spacing}{0.8}
\begin{enumerate}
\item \texthindi{यह प्रमाण पत्र सी एस आई आर-राष्ट्रीय भौतिक प्रयोगशाला, भारत जारी किया गया है जौ कि विज्ञान एवं प्रौद्योगिकी मंत्रालय, भारत सरकार के अधीन वैज्ञानिक व औद्योगिक अनुसंधान परिषद्‌ की संघटक इकाई है एवम्‌ भारत का राष्ट्रीय मापिकी  संस्थान}(NMI) \texthindi{ भी है ।}

\item \texthindi{यह प्रमाण पत्र केवल अंशांकन हेतु जमा किएं गए माषिकी हेतु संदर्थित है।}

\item \texthindi{इस प्रमाण पत्र की प्रतिलिपी, पूर्ण प्रमाण पत्र के अतिरिक्त, तैयार नहीं की जा सकती है, जब तक कि निदेशक, सी एस आई आर-राष्ट्रीय भौतिक प्रयोगशाला, नई दिल्‍ली से अनुमोदित सार के प्रकाशन हेतु लिखित अनुमति प्राप्त नहीं की गयी हो।}

\item \texthindi{उस प्रमाण पत्र में प्रतिवेदित परीक्षण परिणाम केवल मापन की वर्णित परिस्थलियाँ एवं समय हेतु मान्य है।}
\end{enumerate}
\end{spacing}
\centering \includegraphics[width=9cm, height=9cm]{NPL_logo_gray.png}

\begin{center}\LARGE NOTE\end{center}
\begin{spacing}{0.8}
\begin{enumerate}
\item This certificate is issued by CSIR-National Physical Laboratory of India (NPLI) which is a constituent unit of the Council of Scientific \& Industrial Research, the Ministry of Science and Technology, Government of India and is also National Metrology Institute (NMI) of India.
\item This certificate refers only to the particular item (s) submitted for calibration.
\item This certificate shall not be reproduced, except in full, unless written permission for the publication of an approved abstract has been obtained from the Director, CSIR- National Physical Laboratory. New Delhi.
 \item The calibration results reported in this certificate are valid at the time and under the stated conditions of measurement.

\end{enumerate}
\end{spacing}
\end{document}